\chapter{Apéndices}
\begin{verbatim}

    _            __           _ _
   / \   _ __   /_/ _ __   __| (_) ___ ___  ___
  / _ \ | '_ \ / _ \ '_ \ / _` | |/ __/ _ \/ __|
 / ___ \| |_) |  __/ | | | (_| | | (_|  __/\__ \
/_/   \_\ .__/ \___|_| |_|\__,_|_|\___\___||___/
        |_|




Contenido
=========
1. Regla de Simpson.
2. Intervalos de Confianza.
3. Teorema del Límite Central.
4. Tabla de Números Aleatorios.
5. Tabla de la Normal Estándar.
6. Tabla t de Student.
7. Tabla Ji Cuadrada.
8. Tabla Kolmogorov-Smirnov.


 

1. Regla de Simpson.
====================

La importancia de la integración numérica se puede apreciar al observar que frecuentemente el planteamiento de un problema se modela con base en integrales, por lo tanto, para construir una solución es necesario resolverlas. La mayor parte de las integrales no se pueden expresar en términos de funciones elementales por lo que se deben aproximar de manera numérica. 


La integración numérica es un método por medio del cual se obtiene una solución aproximada de una integral definida. En general se busca resolve el problema de:

    _ 
   | )b
   |     f(x) dx  =  I  +  error
 (_|  a



La regla de Simpson es un método de integración numérica relativamente sencillo. Consiste en dividir el intervalo de integración [a,b] en n partes iguales, donde n debe ser un número par, para luego obtener una aproximación del área de cada una de las partes mediante una parábola. A continuación se muestra un gráfico para ilustrar lo dicho, pero la parte superior de cada intervalo es una parábola, no un rectángulo.



|
|                           
|                         . ' .
|                       . __ __ .
|                     .  |  |  |  .
|                   .  __|  |  |__  .
|                 . __|  |  |  |  |__ .
|              . __|  |  |  |  |  |  |__ .
|          .  __|  |  |  |  |  |  |  |  |__  .
|     . __ __|  |  |  |  |  |  |  |  |  |  |__ __ .
|______|__|__|__|__|__|__|__|__|__|__|__|__|__|__|______



A continuación se enuncia la regla de Simpson. Sea [a,b] el intervalo de integración, n un número par que representa el número de partes en que se divide el intervalo y sea:


    b - a
h = -----
      n



Por lo que la integral de f(x)  se puede aproximar mediante:


    _ 
   | )b
   |      f(x) dx  ~  (h/3) * (     1 * f(a + 0h)  
 (_|  a                          +  4 * f(a + 1h)  
                                 +  2 * f(a + 2h)  
                                 +  4 * f(a + 3h) 
                                 +  2 * f(a + 4h) 
                                 ...
                                 +  2 * f(a + (n-2)h) 
                                 +  4 * f(a + (n-1)h) 
                                 +  1 * f(a + (n-0)h)
                               ) 





donde:
      
              /       \     /  h^4   \     / d^4 f(t) \
error = -1 * (  b - a  ) * (  ------  ) * (  --------  )
              \       /     \  180   /     \    dt    /



con "t" elemento de [a,b]


Debe observarse que en general el error es proporcional a la cuarta derivada del intervalo, por ello, se puede satisfacer cualquier condición de error si se escoge un valor de h lo suficientemente pequeño.



() Ejemplo.
-----------

Aplicar la Regla de Simpson para encontrar la integral de la función que se presenta a continuación dentro del intervalo [0,1] y encontrar la magnitud del error.

f(x) = x^2



Como es bien conocido el gráfico de la función está dado por:


 f(x)
   |
  1|                    .
   |                  .
   |               .
   |           .
   |      .
   |.____________________ x
   0                    1



Para resolver el siguiente problema, se utiliza un valor de h=0.25 o lo que es equivalente un valor de n=4. De esta forma el área de integración queda dividida en cuatro partes iguales.



    _ 
   | )b
   |      f(x) dx  =  (h/3) * (      1 * f(a + 0h)  
 (_|  a                           +  4 * f(a + 1h)  
                                  +  2 * f(a + 2h)  
                                  +  4 * f(a + 3h) 
                                  +  2 * f(a + 4h) 
                                  ...
                                  +  2 * f(a + (n-2)h) 
                                  +  4 * f(a + (n-1)h) 
                                  +  1 * f(a + (n-0)h)
                               ) 



    _ 
   | )1
   |      x^2  dx = (0.25/3) * (      1 * f(0 + 0*0.25)  
 (_|  0                            +  4 * f(0 + 1*0.25)  
                                   +  2 * f(0 + 2*0.25)  
                                   +  4 * f(0 + 3*0.25) 
                                   +  1 * f(0 + 4*0.25) 
                                )




    _ 
   | )1   
   |    x^2  dx = (0.25/3) * (      1 * f(0)  
 (_|  0                          +  4 * f(0.25)  
                                 +  2 * f(0.50)  
                                 +  4 * f(0.75) 
                                 +  1 * f(1) 
                              )



    _ 
   | )1                             
   |      x^2  dx = (0.25/3) * (      1 * 0^2  
 (_|  0                            +  4 * 0.25^2          
                                   +  2 * 0.50^2 
                                   +  4 * 0.75^2
                                   +  1 * 1^2
                                )



    _ 
   | )1
   |      x^2  dx = (0.25/3) * (      0  
 (_|  0                            +  0.25  
                                   +  0.50  
                                   +  2.25
                                   +  1 
                                )



    _ 
   | )1
   |      x^2  dx = 0.3333.... 
 (_|  0
       


El valor de 0.3333 es correcto. Como el error depende de la cuarta derivada de la función, en este caso la cuarta derivada de es igual a cero por lo que el error es de cero:


              /       \     /  h^4   \     / d^4 f(t) \
error = -1 * (  b - a  ) * (  ------  ) * (  --------  )
              \       /     \  180   /     \    dt    /



              /       \     / 0.25^4 \     / d^4 f(t) \
error = -1 * (  1 - 0  ) * (  ------  ) * (  --------  )
              \       /     \  180   /     \    dt    /


              /       \     / 0.25^4 \     /   \
error = -1 * (  1 - 0  ) * (  ------  ) * (  0  )
              \       /     \  180   /     \   /



error = 0





2. Intervalos de Confianza.
=========================== 

Muchos de los problemas de simulación no devuelven como resultado un valor único. Si los programas se ejecutan muchas veces este valor tendrá pequeñas oscilaciones. Es por eso que en lugar de presentar estos valores puntuales se recurre a la técnica estadística de intervalos de confianza. 


La estimación de un intervalo describe una gama de posibles valores dentro de los cuales probablemente se encuentre el valor buscado.


Como es sabido la distribución normal cuenta con una serie de cualidades interesantes. Entre ellas la manera como se agrupa su área alrededor de la media.


Es conocido como si se toma una desviación estándar de la media hacia la izquierda y la derecha se agrupa 0.6830 de la distribución. Si se toman dos desviaciones estándar se agrupa el 0.9550 del área de la distribución y se toman 3 desviaciones estándar se agrupa el 0.9970 de la distribución.


Es decir:

 _        _
[x - 1s , x + 1s] agrupa el 0.6830 de la distribución.

 _        _
[x - 2s , x + 2s] agrupa el 0.9550 de la distribución.

 _        _
[x - 3s , x + 3s] agrupa el 0.9970 de la distribución.



|                           
|                         . ' .
|                       .       .
|                     .           .
|                   .               .
|                 .|                 |.
|              .   |                 |   .
|          .|      |                 |      |.
|    |.     |      |                 |      |     .|
|____|______|______|_________________|______|______|________
     |      |      |        _        |      |      |
     |      |      |        x        |      |      |
     |      |    _ |               _ |      |      |
     |      |    x-1s              x+1s     |      |
     |      |      |_________________|      |      |
     |      |            0.6830             |      |
     |    _ |                             _ |      |
     |    x-2s                            x+2s     |
     |      |_______________________________|      |
     |                   0.9550                    |
   _ |                                           _ |
   x-3s                                          x+3s
     |_____________________________________________|
                         0.9970



() Ejemplo.
-----------

En un proceso de simulación se han obtenido los siguientes resultados:

n = 100  tamaño de la muestra.
_
x = 21   media muestral.

s = 6    desviación estándar de la muestra.


Se desea construir un intervalo de confianza del 0.95, o lo que es equivalente del 95% para el resultado encontrado. Para ello debemos realizar los siguientes cálculos.


  _      /    s    \
s(x) =  (  -------  )
         \ sqrt(n) /


  _      /    6      \
s(x) =  (  ---------  )
         \ sqrt(100) /


  _      /    6    \
s(x) =  (  -------  )
         \   10    /


  _ 
s(x) = 0.60



Buscamos en la tabla normal los valores que acumulen un 0.95 del área y denominamos a ese valor z0.


|
|
|                         . ' .
|                       .       .
|                     .           .
|            |      .               .      |
|            |    .                   .    |
|            | .                         . |
|          . |                             | .
|     .      |                             |      .
|____________|_____________________________|________________
             |                             |
           -1.96                          1.96
           -z(0.05/2)                     z(0.05/2)


Posteriormente obtenemos los límites de la distribución, el límite inferior y el límite superior de la distribución.

       _            _
Li  =  x  -  z0 * s(x)
                   

    =  21 - 1.96 * 0.60


    =  19.8240


       _            _
Ls  =  x  +  z0 * s(x)
                   

    =  21 + 1.96 * 0.60


    =  22.1760


Por lo tanto se puede concluir con un 95% de confianza que el valor del promedio se encuentra entre los valores de [19.8240 , 22.1760].

 


3. Teorema del Límite Central.
==============================

A continuación se presenta un ejemplo particular que evidencia el comportamiento del Teorema del Límite Central.


Supongamos que se cuenta con una población de 5 elementos, N = 5, con una distribución uniforme de probabilidad. Es decir, la distribución de probabilidad de la problación se puede representar mediante el siguiente gráfico.



 p(x)
   |
   |
0.4|
   |
0.3|
   |
0.2|   __ __ __ __ __
   |  |  |  |  |  |  |
0.1|  |  |  |  |  |  |
   |__|__|__|__|__|__|__
       1  2  3  4  5




Si se toman muestras de tamaño n=3, el total de muestras posibles está dado por la siguiente tabla.


----------------------
No   Muestra  Promedio
----------------------
 1   1,2,3    2.0000
 2   1,2,4    2.3333
 3   1,2,5    2.6666
 4   1,3,4    2.6666
 5   1,3,5    3.0000
 6   1,4,5    3.3333
 7   2,3,4    3.0000
 8   2,3,5    3.3333
 9   2,4,5    3.6666
10   3,4,5    4.0000
----------------------



Con los datos anteriores se puede construir la siguiente distribución de medias muestrales.


  _
p(x)
  |
  |
 2|                       __ __ __
  |                      |  |  |  |
  |                      |  |  |  |
  |                      |  |  |  |
 1|                 __ __|  |  |  |__ __
  |                |  |  |  |  |  |  |  |
  |                |  |  |  |  |  |  |  |
  |________________|__|__|__|__|__|__|__|_______________
           1        2        3         4        5        



Si se toman muestras de tamaño n=4, el total de muestras posibles está dado por la siguiente tabla.


-----------------------
No   Muestra   Promedio
-----------------------
 1   1,2,3,4   2.5000
 2   1,2,3,5   2.7500
 3   1,2,4,5   3.0000
 4   1,3,4,5   3.2500
 5   2,3,4,5   3.5000
-----------------------



Con los datos anteriores se puede construir la siguiente distribución de medias muestrales.



  _
p(x)
  |
  |
 2|
  |
  |
  |
 1|                      _ _ _ _ _
  |                     | | | | | |
  |                     | | | | | |
  |_____________________|_|_|_|_|_|_____________________
           1        2        3        4        5



Si se toman muestras de tamaño n=5 el total de muestras posibles está dado por la siguiente tabla.


-------------------------
No   Muestra     Promedio
-------------------------
 1   1,2,3,4,5   3.0000
-------------------------



Que se representa como la siguiente distribución que es una normal con una dispersión de cero.


  _
p(x)
  |
  |
 2|
  |
  |
  |
 1|                          _
  |                         | |
  |                         | |
  |_________________________|_|_________________________
           1        2        3        4        5






 
4. Tabla de Números Aleatorios.
===============================


-------------------------------------------------------------------
       (A)    (B)    (C)    (D)    (E)    (F)    (G)    (H)    (I)
-------------------------------------------------------------------
( 1)  80830  17105  90837  82579  25808  52267  53730  95229  66307
( 2)  52587  24081  81638  45959  96192  99321  91023  80682  31865
( 3)  74578  88372  93255  06672  65765  96373  92757  27844  62734
( 4)  14586  96637  27479  97392  07601  74837  57120  77070  15451
( 5)  45201  29302  28134  82576  47326  30064  27699  22906  06730
-------------------------------------------------------------------
( 6)  99544  43197  16259  66529  98513  80960  78567  81980  07337
( 7)  66559  32239  91660  58593  32174  53534  15208  29362  09081
( 8)  28474  01752  19313  52045  68273  17741  27542  23064  16474
( 9)  63121  35728  64555  65607  44858  41258  28725  03533  71540
(10)  58338  98108  01257  20929  96367  39957  03483  78198  56649
-------------------------------------------------------------------
(11)  42264  84551  36805  77532  50069  89911  97290  72029  65361
(12)  59609  21074  31801  74275  14289  35930  08469  53065  79935
(13)  23433  69636  83842  51202  64524  04095  38121  66198  80690
(14)  77976  57209  67179  41068  03161  38141  26394  65120  70722
(15)  94182  86767  59652  50027  08439  36546  67865  30979  74773
-------------------------------------------------------------------
(16)  71985  03308  15141  56983  71731  52338  29370  18946  52292
(17)  12788  13388  32440  59814  42995  67400  70635  91528  80708
(18)  25699  60940  95429  49429  62226  05422  07387  96790  70549
(19)  46096  38992  52456  84872  39939  37072  19905  76670  46365
(20)  12079  23441  53134  87951  61269  40710  60289  82625  81237
-------------------------------------------------------------------
(21)  79171  51101  60506  61871  28820  85390  87235  56228  22231
(22)  41351  39876  05266  32368  34753  24536  41622  02695  72384
(23)  18217  52469  01755  61255  41456  15813  21008  80890  52170
(24)  41179  61656  68643  35999  11446  75002  74341  55369  82327
(25)  52916  08200  89674  84194  54970  63288  61007  50705  61473
-------------------------------------------------------------------
(26)  37209  62796  64814  80587  57974  75221  41578  72623  32782
(27)  57248  96369  32092  49416  98985  14912  36414  63754  37485
(28)  13550  34436  03331  64208  13974  82282  57114  34596  73894
(29)  76967  28445  45444  14807  89797  08142  09392  38087  86123
(30)  39280  89321  78170  86568  54594  93441  24345  74700  49115
-------------------------------------------------------------------


 


5. Tabla de la Normal Estándar.
===============================


|
|                         . ' .
|                       .       .
|         <--.        .           .
|            |      .               .
|            |    .                   .
|            | .                         .
|          . |                               .
|     .      |                                    .
|____________|_____________________________________________
             |
         P(Z<=z0)


--------------------------------------------------------------------------------------
-0.09   -0.08   -0.07   -0.06   -0.05   -0.04   -0.03   -0.02   -0.01   -0.00     z0
--------------------------------------------------------------------------------------
 0.0010  0.0010  0.0011  0.0011  0.0011  0.0012  0.0012  0.0013  0.0013  0.0013  -3.0
 0.0014  0.0014  0.0015  0.0015  0.0016  0.0016  0.0017  0.0018  0.0018  0.0019  -2.9
 0.0019  0.0020  0.0021  0.0021  0.0022  0.0023  0.0023  0.0024  0.0025  0.0026  -2.8
 0.0026  0.0027  0.0028  0.0029  0.0030  0.0031  0.0032  0.0033  0.0034  0.0035  -2.7
 0.0036  0.0037  0.0038  0.0039  0.0040  0.0041  0.0043  0.0044  0.0045  0.0047  -2.6
 0.0048  0.0049  0.0051  0.0052  0.0054  0.0055  0.0057  0.0059  0.0060  0.0062  -2.5
 0.0064  0.0066  0.0068  0.0069  0.0071  0.0073  0.0075  0.0078  0.0080  0.0082  -2.4
 0.0084  0.0087  0.0089  0.0091  0.0094  0.0096  0.0099  0.0102  0.0104  0.0107  -2.3
 0.0110  0.0113  0.0116  0.0119  0.0122  0.0125  0.0129  0.0132  0.0136  0.0139  -2.2
 0.0143  0.0146  0.0150  0.0154  0.0158  0.0162  0.0166  0.0170  0.0174  0.0179  -2.1
 0.0183  0.0188  0.0192  0.0197  0.0202  0.0207  0.0212  0.0217  0.0222  0.0228  -2.0
 0.0233  0.0239  0.0244  0.0250  0.0256  0.0262  0.0268  0.0274  0.0281  0.0287  -1.9
 0.0294  0.0301  0.0307  0.0314  0.0322  0.0329  0.0336  0.0344  0.0351  0.0359  -1.8
 0.0367  0.0375  0.0384  0.0392  0.0401  0.0409  0.0418  0.0427  0.0436  0.0446  -1.7
 0.0455  0.0465  0.0475  0.0485  0.0495  0.0505  0.0516  0.0526  0.0537  0.0548  -1.6
 0.0559  0.0571  0.0582  0.0594  0.0606  0.0618  0.0630  0.0643  0.0655  0.0668  -1.5
 0.0681  0.0694  0.0708  0.0721  0.0735  0.0749  0.0764  0.0778  0.0793  0.0808  -1.4
 0.0823  0.0838  0.0853  0.0869  0.0885  0.0901  0.0918  0.0934  0.0951  0.0968  -1.3
 0.0985  0.1003  0.1020  0.1038  0.1056  0.1075  0.1093  0.1112  0.1131  0.1151  -1.2
 0.1170  0.1190  0.1210  0.1230  0.1251  0.1271  0.1292  0.1314  0.1335  0.1357  -1.1
 0.1379  0.1401  0.1423  0.1446  0.1469  0.1492  0.1515  0.1539  0.1562  0.1587  -1.0
 0.1611  0.1635  0.1660  0.1685  0.1711  0.1736  0.1762  0.1788  0.1814  0.1841  -0.9
 0.1867  0.1894  0.1922  0.1949  0.1977  0.2005  0.2033  0.2061  0.2090  0.2119  -0.8
 0.2148  0.2177  0.2206  0.2236  0.2266  0.2296  0.2327  0.2358  0.2389  0.2420  -0.7
 0.2451  0.2483  0.2514  0.2546  0.2578  0.2611  0.2643  0.2676  0.2709  0.2743  -0.6
 0.2776  0.2810  0.2843  0.2877  0.2912  0.2946  0.2981  0.3015  0.3050  0.3085  -0.5
 0.3121  0.3156  0.3192  0.3228  0.3264  0.3300  0.3336  0.3372  0.3409  0.3446  -0.4
 0.3483  0.3520  0.3557  0.3594  0.3632  0.3669  0.3707  0.3745  0.3783  0.3821  -0.3
 0.3859  0.3897  0.3936  0.3974  0.4013  0.4052  0.4090  0.4129  0.4168  0.4207  -0.2
 0.4247  0.4286  0.4325  0.4364  0.4404  0.4443  0.4483  0.4522  0.4562  0.4602  -0.1
 0.4641  0.4681  0.4721  0.4761  0.4801  0.4840  0.4880  0.4920  0.4960  0.5000   0.0
--------------------------------------------------------------------------------------





|
|                         . ' .
|                       .       .
|                     .           .     <--.
|                   .               .      |
|                 .                   .    |
|              .                         . |
|          .                               | .
|     .                                    |      .
|__________________________________________|________________
                                           |
                                       P(Z<=z0)


------------------------------------------------------------------------------------
 z0   0.00    0.01    0.02    0.03    0.04    0.05    0.06    0.07    0.08    0.09
------------------------------------------------------------------------------------
 0.0  0.5000  0.5040  0.5080  0.5120  0.5160  0.5199  0.5239  0.5279  0.5319  0.5359
 0.1  0.5398  0.5438  0.5478  0.5517  0.5557  0.5596  0.5636  0.5675  0.5714  0.5753
 0.2  0.5793  0.5832  0.5871  0.5910  0.5948  0.5987  0.6026  0.6064  0.6103  0.6141
 0.3  0.6179  0.6217  0.6255  0.6293  0.6331  0.6368  0.6406  0.6443  0.6480  0.6517
 0.4  0.6554  0.6591  0.6628  0.6664  0.6700  0.6736  0.6772  0.6808  0.6844  0.6879
 0.5  0.6915  0.6950  0.6985  0.7019  0.7054  0.7088  0.7123  0.7157  0.7190  0.7224
 0.6  0.7257  0.7291  0.7324  0.7357  0.7389  0.7422  0.7454  0.7486  0.7517  0.7549
 0.7  0.7580  0.7611  0.7642  0.7673  0.7704  0.7734  0.7764  0.7794  0.7823  0.7852
 0.8  0.7881  0.7910  0.7939  0.7967  0.7995  0.8023  0.8051  0.8078  0.8106  0.8133
 0.9  0.8159  0.8186  0.8212  0.8238  0.8264  0.8289  0.8315  0.8340  0.8365  0.8389
 1.0  0.8413  0.8438  0.8461  0.8485  0.8508  0.8531  0.8554  0.8577  0.8599  0.8621
 1.1  0.8643  0.8665  0.8686  0.8708  0.8729  0.8749  0.8770  0.8790  0.8810  0.8830
 1.2  0.8849  0.8869  0.8888  0.8907  0.8925  0.8944  0.8962  0.8980  0.8997  0.9015
 1.3  0.9032  0.9049  0.9066  0.9082  0.9099  0.9115  0.9131  0.9147  0.9162  0.9177
 1.4  0.9192  0.9207  0.9222  0.9236  0.9251  0.9265  0.9279  0.9292  0.9306  0.9319
 1.5  0.9332  0.9345  0.9357  0.9370  0.9382  0.9394  0.9406  0.9418  0.9429  0.9441
 1.6  0.9452  0.9463  0.9474  0.9484  0.9495  0.9505  0.9515  0.9525  0.9535  0.9545
 1.7  0.9554  0.9564  0.9573  0.9582  0.9591  0.9599  0.9608  0.9616  0.9625  0.9633
 1.8  0.9641  0.9649  0.9656  0.9664  0.9671  0.9678  0.9686  0.9693  0.9699  0.9706
 1.9  0.9713  0.9719  0.9726  0.9732  0.9738  0.9744  0.9750  0.9756  0.9761  0.9767
 2.0  0.9772  0.9778  0.9783  0.9788  0.9793  0.9798  0.9803  0.9808  0.9812  0.9817
 2.1  0.9821  0.9826  0.9830  0.9834  0.9838  0.9842  0.9846  0.9850  0.9854  0.9857
 2.2  0.9861  0.9864  0.9868  0.9871  0.9875  0.9878  0.9881  0.9884  0.9887  0.9890
 2.3  0.9893  0.9896  0.9898  0.9901  0.9904  0.9906  0.9909  0.9911  0.9913  0.9916
 2.4  0.9918  0.9920  0.9922  0.9925  0.9927  0.9929  0.9931  0.9932  0.9934  0.9936
 2.5  0.9938  0.9940  0.9941  0.9943  0.9945  0.9946  0.9948  0.9949  0.9951  0.9952
 2.6  0.9953  0.9955  0.9956  0.9957  0.9959  0.9960  0.9961  0.9962  0.9963  0.9964
 2.7  0.9965  0.9966  0.9967  0.9968  0.9969  0.9970  0.9971  0.9972  0.9973  0.9974
 2.8  0.9974  0.9975  0.9976  0.9977  0.9977  0.9978  0.9979  0.9979  0.9980  0.9981
 2.9  0.9981  0.9982  0.9982  0.9983  0.9984  0.9984  0.9985  0.9985  0.9986  0.9986
 3.0  0.9987  0.9987  0.9987  0.9988  0.9988  0.9989  0.9989  0.9989  0.9990  0.9990
------------------------------------------------------------------------------------

 




6. Tabla t de Student.
======================


                          
|                         . ' .
|                       .       .
|         <--.        .           .        .-->
|            |      .               .      |
|            |    .                   .    |
|            | .                         . |
|          . |                             | .
|     .      |                             |      .
|____________|_____________________________|________________
           -t0                             t0
          alfa/2                         alfa/2



--------------------------------------------------
Grados de    ...... Area en Ambos Extremos .......
Libertad     0.10       0.05       0.02       0.01
--------------------------------------------------
    1      6.3137    12.7062    31.8210    63.6559
    2      2.9200     4.3027     6.9645     9.9250
    3      2.3534     3.1824     4.5407     5.8408
    4      2.1318     2.7765     3.7469     4.6041
    5      2.0150     2.5706     3.3649     4.0321
    6      1.9432     2.4469     3.1427     3.7074
    7      1.8946     2.3646     2.9979     3.4995
    8      1.8595     2.3060     2.8965     3.3554
    9      1.8331     2.2622     2.8214     3.2498
   10      1.8125     2.2281     2.7638     3.1693
   11      1.7959     2.2010     2.7181     3.1058
   12      1.7823     2.1788     2.6810     3.0545
   13      1.7709     2.1604     2.6503     3.0123
   14      1.7613     2.1448     2.6245     2.9768
   15      1.7531     2.1315     2.6025     2.9467
   16      1.7459     2.1199     2.5835     2.9208
   17      1.7396     2.1098     2.5669     2.8982
   18      1.7341     2.1009     2.5524     2.8784
   19      1.7291     2.0930     2.5395     2.8609
   20      1.7247     2.0860     2.5280     2.8453
   21      1.7207     2.0796     2.5176     2.8314
   22      1.7171     2.0739     2.5083     2.8188
   23      1.7139     2.0687     2.4999     2.8073
   24      1.7109     2.0639     2.4922     2.7970
   25      1.7081     2.0595     2.4851     2.7874
   26      1.7056     2.0555     2.4786     2.7787
   27      1.7033     2.0518     2.4727     2.7707
   28      1.7011     2.0484     2.4671     2.7633
   29      1.6991     2.0452     2.4620     2.7564
   30      1.6973     2.0423     2.4573     2.7500
   60      1.6706     2.0003     2.3901     2.6603
  120      1.6576     1.9799     2.3578     2.6174
--------------------------------------------------




7. Tabla Ji Cuadrada.
=====================
 

|                 .  .
|               .      .
|             .          .        .-->
|            .             .      |
|           .                .    |
|          .                    . |
|        .                        | .
|     .                           |      .
|_________________________________|______________________
                                  |
                           P(Ji^2>=ji[0])

-------------------------------------------------------------
Grados de    ......... Area en el Extremo Derecho ...........
Libertad     0.99      0.975       0.95      0.90        0.80
-------------------------------------------------------------
    1      0.0002     0.0010     0.0039     0.0158     0.0642
    2      0.0201     0.0506     0.1026     0.2107     0.4463
    3      0.1148     0.2158     0.3518     0.5844     1.0052
    4      0.2971     0.4844     0.7107     1.0636     1.6488
    5      0.5543     0.8312     1.1455     1.6103     2.3425
    6      0.8721     1.2373     1.6354     2.2041     3.0701
    7      1.2390     1.6899     2.1673     2.8331     3.8223
    8      1.6465     2.1797     2.7326     3.4895     4.5936
    9      2.0879     2.7004     3.3251     4.1682     5.3801
   10      2.5582     3.2470     3.9403     4.8652     6.1791
   11      3.0535     3.8157     4.5748     5.5778     6.9887
   12      3.5706     4.4038     5.2260     6.3038     7.8073
   13      4.1069     5.0087     5.8919     7.0415     8.6339
   14      4.6604     5.6287     6.5706     7.7895     9.4673
   15      5.2294     6.2621     7.2609     8.5468    10.3070
   16      5.8122     6.9077     7.9616     9.3122    11.1521
   17      6.4077     7.5642     8.6718    10.0852    12.0023
   18      7.0149     8.2307     9.3904    10.8649    12.8570
   19      7.6327     8.9065    10.1170    11.6509    13.7158
   20      8.2604     9.5908    10.8508    12.4426    14.5784
   21      8.8972    10.2829    11.5913    13.2396    15.4446
   22      9.5425    10.9823    12.3380    14.0415    16.3140
   23     10.1957    11.6885    13.0905    14.8480    17.1865
   24     10.8563    12.4011    13.8484    15.6587    18.0618
   25     11.5240    13.1197    14.6114    16.4734    18.9397
   26     12.1982    13.8439    15.3792    17.2919    19.8202
   27     12.8785    14.5734    16.1514    18.1139    20.7030
   28     13.5647    15.3079    16.9279    18.9392    21.5880
   29     14.2564    16.0471    17.7084    19.7677    22.4751
   30     14.9535    16.7908    18.4927    20.5992    23.3641
-------------------------------------------------------------




|                 .  .
|               .      .
|             .          .        .-->
|            .             .      |
|           .                .    |
|          .                    . |
|        .                        | .
|     .                           |      .
|_________________________________|______________________
                                  |
                           P(Ji**2>=ji0)

-------------------------------------------------------------
Grados de    ......... Area en el Extremo Derecho ...........
Libertad     0.20       0.10       0.05      0.025       0.01
-------------------------------------------------------------
    1      1.6424     2.7055     3.8415     5.0239     6.6349
    2      3.2189     4.6052     5.9915     7.3778     9.2104
    3      4.6416     6.2514     7.8147     9.3484    11.3449
    4      5.9886     7.7794     9.4877    11.1433    13.2767
    5      7.2893     9.2363    11.0705    12.8325    15.0863
    6      8.5581    10.6446    12.5916    14.4494    16.8119
    7      9.8032    12.0170    14.0671    16.0128    18.4753
    8     11.0301    13.3616    15.5073    17.5345    20.0902
    9     12.2421    14.6837    16.9190    19.0228    21.6660
   10     13.4420    15.9872    18.3070    20.4832    23.2093
   11     14.6314    17.2750    19.6752    21.9200    24.7250
   12     15.8120    18.5493    21.0261    23.3367    26.2170
   13     16.9848    19.8119    22.3620    24.7356    27.6882
   14     18.1508    21.0641    23.6848    26.1189    29.1412
   15     19.3107    22.3071    24.9958    27.4884    30.5780
   16     20.4651    23.5418    26.2962    28.8453    31.9999
   17     21.6146    24.7690    27.5871    30.1910    33.4087
   18     22.7595    25.9894    28.8693    31.5264    34.8052
   19     23.9004    27.2036    30.1435    32.8523    36.1908
   20     25.0375    28.4120    31.4104    34.1696    37.5663
   21     26.1711    29.6151    32.6706    35.4789    38.9322
   22     27.3015    30.8133    33.9245    36.7807    40.2894
   23     28.4288    32.0069    35.1725    38.0756    41.6383
   24     29.5533    33.1962    36.4150    39.3641    42.9798
   25     30.6752    34.3816    37.6525    40.6465    44.3140
   26     31.7946    35.5632    38.8851    41.9231    45.6416
   27     32.9117    36.7412    40.1133    43.1945    46.9629
   28     34.0266    37.9159    41.3372    44.4608    48.2783
   29     35.1394    39.0875    42.5569    45.7223    49.5878
   30     36.2502    40.2560    43.7730    46.9792    50.8923
-------------------------------------------------------------





8. Tabla Kolmogorov-Smirnov.
============================


------------------------------------------------------
Tamaño de   ........ Nivel de Significancia ..........
muestra n   0.20     0.15     0.10     0.05     0.01
------------------------------------------------------
    1       0.900    0.925    0.950    0.975    0.995
    2       0.684    0.726    0.776    0.842    0.929
    3       0.565    0.597    0.642    0.708    0.828
    4       0.494    0.525    0.564    0.624    0.733
    5       0.446    0.474    0.510    0.565    0.669

    6       0.410    0.436    0.470    0.521    0.618
    7       0.381    0.405    0.438    0.486    0.577
    8       0.358    0.381    0.411    0.457    0.543
    9       0.339    0.360    0.388    0.432    0.514
   10       0.322    0.342    0.368    0.410    0.490

   11       0.307    0.326    0.352    0.391    0.468
   12       0.295    0.313    0.338    0.375    0.450
   13       0.284    0.302    0.325    0.361    0.433
   14       0.274    0.292    0.314    0.349    0.418
   15       0.266    0.283    0.304    0.338    0.404

   16       0.258    0.274    0.295    0.328    0.392
   17       0.250    0.266    0.286    0.318    0.381
   18       0.244    0.259    0.278    0.309    0.371
   19       0.237    0.252    0.272    0.301    0.363
   20       0.231    0.246    0.264    0.294    0.356

   25       0.210    0.220    0.240    0.270    0.320
   30       0.190    0.200    0.220    0.240    0.290
   35       0.180    0.190    0.210    0.230    0.270

            1.07     1.14     1.22     1.36     1.63
 n>35       ____    _____    _____    _____    _____
              __       __       __       __       __
            \/n      \/n      \/n      \/n      \/n  

------------------------------------------------------




\end{verbatim}